\section{System Design \& Task Setup}
\label{sec:system_design}

\subsection{Hardware}

Experiments were run on a da Vinci Research Kit (dVRK)~\cite{dvrk}. The
dVRK pairs a da Vinci S surgeon-side console with da Vinci Si
patient-side arms, instrumented with bipolar fenestrated forceps. Two
patient-side manipulators, PSM1 and PSM2, are available for control;
trials can use either arm alone or both together.

\subsection{Task}
\label{subsec:task}

The task is a cylinder-on-peg transfer task. The task board holds four
blue and four white cylinders at rest on four center pegs. Eight pegs
ring the outside of the board; each is embedded with an LED that lights
either blue or white to indicate the target cylinder and its destination
peg for the next move (Figure~\ref{fig:task_board}).

The board is instrumented with conductive paint, wired to detect cylinder
placement and removal directly: contact between a cylinder and a peg
closes a circuit specific to that peg. This gives a ground-truth signal
for cylinder state -- seated on a peg, or in transit -- independent of
the vision pipeline.

\begin{figure}[h]
\centering
% \includegraphics[width=0.6\textwidth]{figures/task_board.png}
\caption{Task board layout. Four center pegs hold the cylinders at rest.
Eight outer pegs illuminate blue or white via embedded LEDs to indicate
the target cylinder and destination for the next move.}
\label{fig:task_board}
\end{figure}

\subsection{Vision}
\label{subsec:vision}

Two camera streams feed the pipeline. The dVRK endoscopic camera provides
the stereo view used at the surgeon console. An Intel RealSense D405
depth camera, mounted next to the task board, provides a fixed external
RGB-D view of the workspace, independent of console head motion.
% TODO: auxiliary camera view -- extra pedal toggles a second view
% (top or side of task pad) during robot control, for added depth
% perception. Describe pedal mapping and where the view is displayed.

\subsection{Sensing}
\label{subsec:sensing}

Three further modalities complement the two camera streams. Kinematic
data -- joint positions, velocities, and tool-tip pose -- is read
directly from the patient-side manipulators (PSMs). Task start and end
events are detected by an Arduino-based trigger circuit, independent of
the vision and kinematic streams. Contact forces are measured by an ATI
force sensor mounted underneath the task pad, capturing forces
transmitted through the board as cylinders are placed and removed.
